\zhappendix

\section{从 $X(\Fq)$ 上的函数到 $X(\CC)$ 上的 D-模}
\label{s:frmfun}

\subsection{概览}

\ssszh 本附录是第 \ref{s:bbloc} 节的补充，目的是解释：从旗簇的 $\Fq$-点上的函数过渡到旗簇的 $\CC$-点上的
D-模，何以是一件自然的事。

简短的回答是：人们把后者看作前者某种意义上的范畴化。任何解释都必然涉及几何表示论中三个重要的
主题——D-模、函数-层对应，以及 Riemann--Hilbert 对应。

我们希望下面关于这些主题的讲述，能为非专家提供足够的背景与直觉，以便欣赏正文中的讨论。
文献指引我们留到每一小节的末尾。

\subsection{代数 D-模}

\ssszh 本小节讨论 D-模的基础。先从仿射情形开始。

设 $X$ 是域 $k$ 上的光滑仿射代数簇，为简单起见假定 $k$ 的特征为零。粗略地说，$X$ 上的 D-模就是
这样一个对象：在其上既可以乘以 $X$ 上的函数，也可以沿 $X$ 上的向量场求导。

更形式地说，记 $O_X$ 为 $X$ 上正则函数构成的代数，$T_X$ 为 $X$ 上正则向量场构成的 $O_X$-模。
$X$ 上的 D-模 $M$ 就是带有平坦联络的 $O_X$-模。也就是说，$M$ 是向量空间，带有函数代数的作用
\[
O_X\otimes M\to M,\quad f\otimes m\mapsto f\cdot m,
\]
以及向量场李代数的作用
\[
T_X\otimes M\to M,\quad \xi\otimes m\mapsto \xi\cdot m.
\]
上面是李代数作用，这恰恰就是平坦性（即曲率消失）的假设。还要求这两种运算与「以函数缩放向量场」
相容
\[
f\cdot(\xi\cdot m)=(f\cdot\xi)\cdot m,
\]
并与「用向量场求导函数」相容，即满足「莱布尼茨法则」
\[
\xi\cdot(f\cdot m)=(\xi\cdot f)\cdot m+f\cdot(\xi\cdot m),
\]
其中 $\xi\cdot f$ 表示 $f$ 沿 $\xi$ 的导数。

等价地，可以构造 $X$ 上多项式线性微分算子构成的代数 $D_X$，即 $O_X$ 的全部自同态中由「乘以函数」
与「沿向量场求导」生成的子代数。有了它，D-模就不过是左 $D_X$-模。

\begin{ex} 对仿射空间 $X=\mathbb{A}^n_k$，正则函数与向量场形如
\[
O_X\simeq k[x_1,\ldots,x_n]\quad\text{和}\quad T_X\simeq
\underset{i}{\oplus}\ O_X\cdot\partial_i,
\]
其中 $\partial_i$（$1\leqslant i\leqslant n$）记与坐标 $x_i$ 相伴的标准坐标向量场。此时，$D_X$
作为 $k$-代数由 $x_i,\partial_j$（$1\leqslant i,j\leqslant n$）生成，满足标准的交换关系
\[
[x_i,x_j]=[\partial_i,\partial_j]=0,\qquad
[\partial_i,x_j]=\begin{cases}1 & \text{若 } i=j,\\ 0 & \text{否则}.\end{cases}
\]
具体地说，$D_X$ 以有序单项式
\[
x_1^{a_1}\cdots x_n^{a_n}\cdot\partial_1^{b_1}\cdots\partial_n^{b_n},\qquad
\text{其中 } a_i,b_j\in\mathbb{Z}^{\geqslant 0},
\]
为基，即由系数为多项式的线性微分算子构成。
\end{ex}

记 $X$ 上 D-模构成的阿贝尔范畴为 $\on{D-mod}(X)^\heartsuit$，其无界导出范畴记为
$\on{D-mod}(X)$。下面讨论这两个范畴中三个基本对象，也就是 D-模的三个例子。

\begin{ex} 微分算子代数 $D_X$ 自然地是自身的左模，即一个 D-模。
\end{ex}

下面转向一个不那么平凡、「整体」的例子。

\begin{ex} 对任意 $X$，正则函数环 $O_X$ 带有显然的 D-模结构：按标准方式乘函数、对函数求导。

该模有如下生成元与关系的表示。作用在常函数 $1$ 上给出 D-模的同构
\[
D_X/D_X\cdot T_X\simeq O_X.
\]
特别地，给出 D-模的映射 $O_X\to M$，等同于指定 $M$ 的一个平坦截面，即被 $T_X$ 零化的元素 $m$。

更一般地，任何带平坦联络 $\nabla$ 的向量丛 $E$ 都是 D-模；给出 D-模的映射 $E\to M$，等同于指定
张量积 $M\otimes_{O_X}E^\vee$（带其自然平坦联络）的一个平坦截面。
\end{ex}

最后，我们想描述一类「局部」的例子。

\begin{ex} 对 $X$ 的任意闭点 $x$，有一个相应的 delta 函数 D-模。记 $\fm_x\subset O_X$ 为在 $x$
处消失的函数构成的极大理想，令
\[
\delta'_x:=D_X/(D_X\cdot\fm_x).
\]
特别地，给出 D-模的映射 $\delta'_x\to M$，等同于指定 $M$ 的一个被 $\fm_x$ 零化的截面，
即在 $x$ 处概形论地支撑的截面。

D-模 $\delta'_x$ 是「支撑在 $x$ 处的分布构成的向量空间」的代数几何版本。\footnote{若改用相应的
右 D-模，这一类比会更贴切，但读者可以略过这一点。} 与此相关，若 $\xi_1,\ldots,\xi_n$ 是
$T_X$ 中的一族向量场，它们在 $x$ 处切空间 $T_xX$ 中的像构成一组基，则单项式
\begin{equation} \label{PBW}
\xi_1^{a_1}\cdots\xi_n^{a_n},\qquad\text{其中 } a_i\in\mathbb{Z}^{\geqslant 0}
\end{equation}
构成 $\delta'_x$ 的一组基。人们把 $1$ 在 $D_X$ 中的像想成通常的 delta 分布，整个模便由它的各阶
导数张成。

这里有一个稍微技术性的点，读者初读时可以放心略过。$\delta'_x$ 的底向量空间若用一条线扭一下会
更方便，即 $x$ 处切空间的行列式线：
\[
\delta_x:=\delta'_x\underset{k}{\otimes}\det(T_xX).
\]
理由如下。给定光滑仿射簇的映射 $f:X\to Y$，有相应的 D-模推前
$$f_*:\on{D-mod}(X)\to\on{D-mod}(Y),$$ 参见下文第 \ref{ss:pushpull} 节。把这一扭纳入之后，
便有典范同构
\[
f_*(\delta_x)\simeq\delta_{f(x)},\qquad\text{对 } x\in X.
\]
也就是说，delta 函数推前为 delta 函数，这与分布上的类比所预期的一致。
\end{ex}

\begin{re} \label{r:hwdelta} 这里给出 $\delta_x$ 的另一种构造，它有助于说明行列式线为何出现。

考虑在 $x$ 处消失的向量场构成的李代数，即 $\fm_x\cdot T_X\subset T_X$。它自然地作用在商
$T_X/\fm_x\cdot T_X\simeq T_xX$ 上。具体地说，这一作用通过
\[
\fm_x\cdot T_X/\fm_x^2\cdot T_X\simeq T^*_xX\otimes T_xX\simeq\on{End}(T^*_xX)
\]
分解，而它在 $T_xX$ 上的作用，对任意向量空间 $W$，就是自同态李代数 $\on{End}(W)$ 在对偶空间
$W^*$ 上的自然作用。

考虑 $\fm_x\cdot T_X$ 在行列式线 $\det(T_xX)$ 上的诱导作用。这条线还带有 $\OO_X$ 的自然作用，
这两个作用定义了由 $O_X$ 与 $\fm_x\cdot T_X$ 生成的 $D_X$ 的子代数 $D_{X,x}$ 的作用。要点在于：
若不先过渡到行列式，所需的种种相容性就不再成立。

最后，诱导出的 D-模带有典范同构
\[
D_X\underset{D_{X,x}}{\otimes}\det(T_xX)\simeq\delta_x,
\]
它在所生成的行列式线上作用为恒等，也就是我们已描述了 $D_{X,x}$ 在 delta 分布本身所张成的空间
上的作用。
\end{re}

正如读者可以想见，许多 D-模的行为介于上面最后两个例子之间：它们沿某个子簇像代数联络，而沿其
法向像 delta D-模。

\ssszh \label{ss:pushpull} 给定光滑仿射簇的映射 $X\to Y$，其 D-模范畴之间有自然的拉回与推前
函子，即
\[
f^!:\Dmod(Y)\to\Dmod(X)\qquad\text{和}\qquad f_*:\Dmod(X)\to\Dmod(Y).
\]
关于这两个函子的性质，就本书的目的而言，只需记住下面这一点。流形 $M$ 上有光滑函数代数及其对偶
的分布空间。给定映射 $M\to N$，可以拉回函数，并借助对偶推前分布。

直观地说，$f^!$ 记录的是：$Y$ 上满足某个线性偏微分方程的函数，其拉回满足哪些线性偏微分方程。
类似地，$f_*$ 记录的是：$X$ 上满足某个线性偏微分方程的分布，其推前满足哪些线性偏微分方程。
特别地，$f^!$ 与 $f_*$ 并不总是伴随函子，而是互为对偶。这可以用第 \ref{ss:dualizable} 节的
dg 范畴对偶性精确表述。

为方便读者，下面更详细地讨论 $f^!$ 与 $f_*$，不过读者可以放心略过。拉回 $f^!$ 就是把拟凝聚
拉回再作一个上同调平移，即
\[
f^!(M)\simeq O_X\underset{O_Y}{\otimes}M[d_X-d_Y],
\]
其中 $d_X,d_Y$ 分别记 $X,Y$ 的维数。这里 $O_X$ 以显然的方式作用，$T_X$ 则通过链法则作用，
即经由映射 $$T_X\to O_X\underset{O_Y}{\otimes}T_Y$$ 以及后者在 $f^!(M)$ 上的自然作用。要描述
$f_*$，注意上面拉回的公式可以改写为用双模
\[
O_X\underset{O_Y}{\otimes}D_Y
\]
在 $D_Y$ 上作张量。同一个双模也定义推前函子 $f_*$。更仔细地说，它给出右 D-模范畴之间的函子
$$f_{\on{r},*}:\on{mod-D}(X)\to\on{mod-D}(Y).$$ 这一差别由「左 D-模与右 D-模典范等同」来弥合。
即在光滑仿射簇 $Z$ 上，典范丛 $K_Z$ 带有 $D_Z$ 的自然右作用。粗略地说，人们把典范丛看作分布的
代数替代，而该作用就是分部积分。用它与左 D-模作张量，便定义了所需的等价
\[
-\underset{O_Z}{\otimes}K_Z:\Dmod(Z)\simeq\on{mod-D}(Z),
\]
在两边施用，便得到推前
\[
f_*:\Dmod(X)\simeq\on{mod-D}(X)\to\on{mod-D}(Y)\simeq\Dmod(Y).
\]

\ssszh 仿射情形下 D-模的定义可以直接整体化，如下。设 $X$ 是 $k$ 上的光滑代数簇，不必仿射。

$X$ 上的 D-模是带平坦联络的拟凝聚层 $\sM$。也就是说，记 $\OO_X$ 为正则函数层、$\sT_X$
为正则向量场层、$\underline{k}$ 为 $k$ 对应的常层，则 $\sM$ 是 $k$ 向量空间的层，带有映射
\[
\OO_X\underset{\ck}{\otimes}\sM\to\sM\qquad\text{与}\qquad
\sT_X\underset{\ck}{\otimes}\sM\to\sM,
\]
满足与上一小节相同的相容性。所谓拟凝聚，是进一步要求：对任意仿射开集 $U$ 及 $U$ 上正则函数
$f$（其非零轨迹记 $U_f$），自然映射
\[
\Gamma(U,\sM)\underset{O_U}{\otimes}O_{U_f}\xrightarrow{\sim}\Gamma(U_f,\sM)
\]
为同构。由这一假设，若 $X$ 仿射，则当下这个层论版本的 D-模范畴与上一小节的范畴经由「取整体
截面」函子典范等价。

于是，粗略地说，一般光滑簇上的 D-模范畴由其上各仿射开子簇的 D-模范畴粘合而成。类似地，
D-模 $\sM$ 大致就是这样一个层：其局部截面既可以乘以正则函数，也可以沿正则向量场求导。

上一小节的例子如下搬到这一设定中。

\begin{ex} 有 $X$ 上微分算子的整体层 $\sD_X$，D-模就是 $\sD_X$-模构成的拟凝聚层。
\end{ex}

\begin{ex} 对任意 $X$，正则函数层 $\OO_X$ 自然地是 D-模；任何带平坦联络的向量丛
$(\sE,\nabla)$ 的截面层亦然。
\end{ex}

\begin{ex} 对 $X$ 的闭点 $x$，同样有支撑在 $x$ 处的 delta 函数 D-模，记作 $\delta_x$。
\end{ex}

\label{ss:dmodglob}
\ssszh 最后指出：如第 \ref{ss:dmodglob} 节所讨论，$X$ 上的 D-模自然地构成一个阿贝尔范畴，
使得忘掉 D-模结构、只看阿贝尔群层的函子是 $t$-正合的。也就是说，核、余核与直和按层的通常方式
计算，它们自然地继承 D-模结构，并满足正确的泛性质。

\begin{defn} 记 $\on{D-mod}(X)^\heartsuit$ 为 $X$ 上 D-模构成的阿贝尔范畴，其无界导出范畴记为
$\on{D-mod}(X)$。
\end{defn}

给定映射 $f:X\to Y$，同样有推前与拉回函子
\[ f^!:\Dmod(Y)\to\Dmod(X)\qquad\text{和}\qquad f_*:\Dmod(X)\to\Dmod(Y), \]
定义同第 \ref{ss:pushpull} 节。应当强调：在非仿射情形，$f_*$ 是「用双模
$\sO_X\otimes_{\sO_Y}\sD_Y$ 作张量」的左导出函子与「沿 $f$ 推前层」的右导出函子的复合。因此
它不是两个 D-模阿贝尔范畴之间某个函子的导出函子，必须从一开始就在导出范畴中工作。

\ssszh 关于 D-模的进一步阅读，读者可以参阅下面这些（在众多可能的文献中，按对读者的要求大致
递增排列）：\cite{coutinho}、\cite{arkhipovnotes}、\cite{bernsteinnotes}、
\cite{kashiwarabookd-modules}、\cite{gaitsgorynotes}、\cite{hottatakeuchitanisaki}。
